% !TeX root = ../main.tex
\chapter{Case Studies}
\label{chap:case-studies}

\textbf{Draft outline.} This chapter will adapt the cases already documented
in the project repository. Each case will describe its context, evidence, the
three CD4AI stages, and the connections and limitations of the feedback loop.

% Existing analyses: ../tcc-docs/cases/<case-name>/README.md.
% Source snapshots and revision indexes: each case's assets/sources/README.md.
% Figure provenance and adaptation details: each case's assets/README.md.
% Add bibliography entries and citations when adapting the case narratives.

\section{How to Read the Cases}
\label{sec:cases-overview}

\begin{itemize}
  \item Briefly recall the analysis procedure from Chapter~\ref{chap:methodology}.
  \item Distinguish source inspection, documented capabilities, illustrative
    screenshots, and executed observations. Record source dates and revisions.
  \item Explain that mapping a workflow to CD4AI does not imply that its
    maintainers explicitly adopted the pattern.
\end{itemize}

\subsection{Questions for Examining the Proposed Refinement}
\label{sec:cases-evaluation-purpose}

The next review of DeepEval/Confident AI, Braintrust, and MLflow will examine
how their documented workflows support the distinction proposed in
Section~\ref{sec:cd4ai-result-interpretation}. The analysis should separate
evidence for the feedback loop from evidence for interpreting failures and
making delivery decisions.

\begin{itemize}
  \item What is retained after review: an interaction, a confirmed defect,
    or an executable evaluation case with explicit criteria?
  \item How are expected behavior, unresolved limitations, and previously
    established behavior represented?
  \item How does an evaluation result distinguish a persistent limitation
    from deterioration relative to an accepted reference?
  \item What happens to a failing case: does it block delivery, contribute to
    a comparative score, or remain a reported limitation?
  \item Is there evidence that a case changes purpose after an improvement
    is confirmed, and who makes that decision?
  \item Which choices are configurable mechanisms, such as dataset partitions,
    baselines, thresholds, and reporting, rather than requirements of CD4AI?
\end{itemize}

For each question, record the source and distinguish documented mechanisms
from observed team practices and our interpretation. Where the available
material does not answer a question, report it as not documented rather than
assuming the behavior is either present or absent. In particular, support for
multiple datasets does not establish that a team separates capability and
regression evaluations. Preserve the dated source snapshots and identify
newer documentation separately when extending the analyses.

\section{Divination}
\label{sec:case-divination}
% Source: ../tcc-docs/cases/divination/README.md.

\subsection{Context and Evidence}
\noindent\emph{To develop:} Introduce the Dungeons \& Dragons RAG assistant,
the inspected implementation revision, and the scope of the local work.
Identify screenshots that use synthetic demonstration data.

\subsection{Testing, Monitoring, and Curation}
\noindent\emph{To develop:} Map the evaluation suites, stored interactions,
candidate signals, and human review to the three stages.

\subsection{Connections and Limitations}
\noindent\emph{To develop:} Explain the documented manual path from a reviewed
defect to a versioned evaluation case. Distinguish available implementation
and instructions from any cycle actually executed during the project.

\section{DeepEval and Confident AI}
\label{sec:case-deepeval}
% Source: ../tcc-docs/cases/deepeval/README.md.

\subsection{Context and Evidence}
\noindent\emph{To develop:} Distinguish the evaluation framework from the
associated platform and identify the documentation and source files analyzed.

\subsection{Testing, Monitoring, and Curation}
\noindent\emph{To develop:} Explain the documented development--production
workflow and its return path from failing conversations to evaluation datasets.

\subsection{Connections and Limitations}
\noindent\emph{To develop:} Explain the interpretation of curation in that
return path and the workflow a team would need to assemble. State that the
existing analysis did not execute the cycle on the platform.

\section{Braintrust}
\label{sec:case-braintrust}
% Source: ../tcc-docs/cases/braintrust/README.md.

\subsection{Context and Evidence}
\noindent\emph{To develop:} Introduce the documented workflow and the provenance
of its published interface examples.

\subsection{Testing, Monitoring, and Curation}
\noindent\emph{To develop:} Map experiments, logs, human annotation, and dataset
creation to the CD4AI stages, highlighting the explicit review activity.

\subsection{Connections and Limitations}
\noindent\emph{To develop:} Describe how selected traces and expected behavior
can become dataset entries. Clarify that the published screenshots show
separate examples, rather than one cycle executed for this study.

\section{MLflow}
\label{sec:case-mlflow}
% Source: ../tcc-docs/cases/mlflow/README.md.

\subsection{Context and Evidence}
\noindent\emph{To develop:} Introduce the documented regression-testing workflow
and identify the dated sources used in the analysis.

\subsection{Testing, Monitoring, and Curation}
\noindent\emph{To develop:} Map tracing, feedback, review queues, evaluation
datasets, and regression tests to the CD4AI cycle.

\subsection{Connections and Limitations}
\noindent\emph{To develop:} Explain the return path from an observed failure to
a retained input and expected behavior. State that the existing analysis
examines documentation and did not run an MLflow server or evaluation.

% Cross-case synthesis and implications for pattern refinement belong in
% Chapter 6 (Discussion), after the individual case narratives are developed.
